\documentclass[11pt]{article}
\usepackage[T1]{fontenc}
\usepackage[utf8]{inputenc}
\usepackage{lmodern,microtype,amsmath,amssymb,amsthm,mathtools,booktabs}
\usepackage[a4paper,margin=24mm]{geometry}
\usepackage[colorlinks=true,linkcolor=blue,urlcolor=blue]{hyperref}
\usepackage{enumitem}
\setlist{nosep,leftmargin=*}
\newtheorem{proposition}{Proposition}
\newcommand{\pos}[1]{(#1)_+}
\title{Erd\H{o}s Problem 708: Execution of the Cold-Quotient Route}
\author{Research continuation prepared for Denis Paliy}
\date{24 November 2026}
\begin{document}
\maketitle
\begin{abstract}
The proposed general descent inequality remains unproved. This continuation
checks three attempted ways to establish it and gives exact obstructions to
their intermediate steps. A real full tail can have large source excess and
no positive cold quotient, while its target banks compensate the excess.
Unions of hot roots below the square-root scale can have discrepancy much
larger than that scale. An abstract model has a unique global optimum with
all cold quotients negative, despite exact rows and divisor domination; its
failure of the numerical floor laws is explicit. These results neither prove
nor disprove the general asymptotic or exact $2n$.
\end{abstract}

\section{The step attempted}
Fix $c=2+\delta>2$, an actual positive $M$-term interval $I$, and
\[
 X(t)=\sum_{p\le M}w_pv_p(t),\quad0\le w_p\le1,
 \quad A=\{a\le M:X(a)>c\}.
\]
Write $S_c=\sum_{a\le M}(X(a)-c)_+$,
$H=\sum_{b\in I}(X(b)-1)_+$ and $J_c=S_c-H$.
For a cold $2\le d\le M$, meaning $s=X(d)<1$, put $\tau=1-s$,
$N_d=\lfloor M/d\rfloor$. The quotients of the interval multiples of $d$
form $N_d$ or $N_d+1$ consecutive integers. Retain $N_d$ of them as $I_d$,
and use prices
\[
 w_p^{(d)}=\min(w_p/\tau,1)\ (p\le N_d),\qquad w_p^{(d)}=0\ (p>N_d).
\]
Let $Q_d=J_c(w^{(d)};N_d,I_d)$, with the better endpoint deletion if
needed. Then $Q_d\le\Delta_c(N_d)$ and $N_d\le M/2$, where $\Delta_c$
is the worst genuine surplus. The unproved target is
\[
 \boxed{J_c\le\frac43\max\left(0,\max_{d\text{ cold}}Q_d\right)
                      +C_c\sqrt M.}\tag{D}
\]
It is required on genuine global maxima with their own full tails.
The previously verified recursion and threshold-shift arguments imply
that (D), combined with the prior near-quadratic exclusion, would close
the asymptotic. That implication is not a proof of (D).

The exact quotient identity remains valid:
\[
 \tau Q_d=K_d+R_d^{\rm tail}+R_d^{\rm end}+R_d^{\rm clip},\qquad R_d^*\ge0,
\]
where
\[
 K_d=\sum_{\substack{a\in A\\d\mid a}}(X(a)-c)
 +(c-1)X(d)\#\{a\in A:d\mid a\}
 -\sum_{\substack{b\in I\\d\mid b}}(X(b)-1)_+.
\]
The reserves are new quotient-tail mass, a removed endpoint penalty, and
clipping gain. No original coordinate balances are transferred to a quotient.

\section{A real full tail need not localize its source excess}
\begin{proposition}
For fixed $2<c<3$, there are actual raw full tails with
$S_c/\sqrt M\to\infty$, but $Q_d\le0$ for every cold $d\ge2$ and every
actual target interval. Their weights fail genuine optimality: target banks
pay more than the entire source excess.
\end{proposition}
\begin{proof}
Let $P\ge1$ be an integer, let $R\subset(4P,5P]$ be a nonempty set of
primes, put $k=|R|$, $M=125P^3$, and give precisely these primes unit
prices. Every cubic monomial in $R$, repetitions allowed, lies in
$(64P^3,125P^3]\subset(M/2,M]$. Every quartic monomial exceeds $M$.
There is no extra inactive cofactor for a cubic monomial, since even
multiplication by two exceeds $M$. Thus the entire full tail is the
set of cubic monomials, with
\[
 n=\binom{k+2}{3},\qquad S_c=(3-c)\binom{k+2}{3},\qquad
 D_p(A)=\frac{(k+1)(k+2)}2\quad(p\in R).\tag{1}
\]
A cold integer contains only zero-price primes. Its normalized positive
prices remain one, possibly with coordinates deleted. Its source prefix
has length at most $M/2$, hence contains no triple of active occurrences.
The new source hinge is zero, so $Q_d\le0$.
Taking all primes in $(4P,5P]$, the prime number theorem gives
$k\sim P/\log P$, and $S_c\asymp_c P^3/(\log P)^3\gg\sqrt M$.
\end{proof}

\paragraph{Exact compensation and the violated optimality row.}
For fixed $p\in R$, take the one common union of banks $pq$, $q\in R$,
including $q=p$, and denote its size by $U_p$. Each bank is strictly hot
and has at least $5P$ positions. The least common multiple of two distinct
banks exceeds $64P^3$, so their intersection has at most two positions.
Consequently
\[
 U_p\ge5Pk-k(k-1)>\frac{(k+1)(k+2)}2=D_p(A),\tag{2}
\]
since $k\le P$; the difference is at least
$\tfrac72k^2-\tfrac12k-1>0$.
If these prices were optimal, every position in the union would have
$z_b=1$, already exceeding its active $p$-row demand. Thus exact
complementary slackness excludes these prices.

There is also direct common-budget payment. At a hot target $b$, the
number $r(b)$ of prime-row unions containing it is at most $X(b)$.
Since $X(b)\ge2$, $X(b)-1\ge r(b)/2$. Summing gives
\[
 H\ge\frac12\sum_{p\in R}U_p
 \ge\frac{k}{2}[5Pk-k(k-1)]
 \ge2k^3+\frac{k^2}{2}>\binom{k+2}{3}>S_c.\tag{3}
\]
Therefore $J_c<0$ for every real interval. This does not refute (D).
It disproves the source-only step
$S_c\le C\sqrt M+C'\max(0,\max_dQ_d)$, and shows why target payment
must remain coupled to localization.

\section{Hot roots below $\sqrt M$ can have much larger discrepancy}
\begin{proposition}
A generic prefix-to-interval union-transfer error $O(\sqrt M)$ is false
even for strictly hot roots at most $\sqrt M$, with a nonempty full tail
and no unused active prime.
\end{proposition}
\begin{proof}
Take $c=9/4$, $M=4z^2$, $N=2z^2$, and
$w_2=3/4$, $w_p=1$ for odd primes $p\le z$, all other prices zero.
The roots $\mathcal H=\{2p:p\le z\}$ satisfy $h\le\sqrt M$ and are hot:
$X(4)=3/2$, $X(2p)=7/4$ for odd $p$.
Every active odd prime occurs in the full tail via $2p^2$, of load $11/4$;
the prime two occurs too.

Let $P(z)=\prod_{p\le z}p$ and let $R_z(x,N)$ count integers coprime to
$P(z)$ in $\{x+1,\ldots,x+N\}$. For the actual interval
$I_x=\{2x+1,\ldots,2x+M\}$,
\[
 \left|\bigcup_{h\in\mathcal H}(I_x\cap h\mathbb Z)\right|=N-R_z(x,N).
\]
The exact CRT average is
$\operatorname{Avg}_{x\bmod P(z)}R_z(x,N)=N\prod_{p\le z}(1-1/p)$.
By Mertens' product theorem this is
$(e^{-\gamma}+o(1))N/\log z$.
A $z$-rough integer at most $2z^2$ is one, a prime exceeding $z$, or a
product of two such primes. The last type has only
$O(z^2/\log^2z)$ elements. Hence, by the prime number theorem,
$R_z(0,N)=(1/2+o(1))N/\log z$.
Some actual shift therefore has prefix-union minus target-union size
at least
\[
 \left(e^{-\gamma}-\frac12+o(1)\right)\frac N{\log z}\gg\sqrt M.\tag{4}
\]
\end{proof}
At $z=1000$, the exact prefix rough count is $152{,}528$ and the ceiling
of the exact rational CRT mean is $161{,}931$. Thus an actual deficit
of at least $9403$ exists, while $\sqrt M=2000$. A standard-library
Python verifier accompanies this existence certificate.

These prices are not asserted to be positive global maximizers. Further,
the actual eligible sources with a hot root and a cold complement have
forms $2pqv$ or $4pv$, with odd active $p,q$ and inactive $v$. Their
number is $o(M)$, whereas the bank union contains all $M/4$ multiples
of four. Replacing actual source demand by the whole prefix union loses
an essential reserve. This proposition does not refute (D).

\section{An abstract unique global optimum defeats formal descent}
This section is expressly nonarithmetic. There are six formal prime
coordinates and 90 source profiles:
\begin{itemize}
\item three copies of $(1,1,1,1,1,1)$;
\item one copy of $2\mathbf1_S$ for each four-element subset $S$;
\item twelve copies of each singleton $e_i$.
\end{itemize}
The 90 target profiles are six copies of $(2,2,2,2,2,2)$, four copies of
each $2e_i$, three copies of each $e_i+e_j$, and fifteen zero profiles.
Set $c=21/10$. At $w_i=1/2$, the strict full tail has 18 sources, load
sum 69 and demand 23 in every coordinate. Select all six large targets,
all 24 doubled singletons and mass $1/5$ at each pair target. Every row
supplies $12+8+3=23$, at total cost $F=39$. Therefore
\[
 J_{21/10}=39-\frac{21}{10}\,18=\frac65.\tag{5}
\]
Every divisor orthant is dominated on the target; all full raw coordinate
totals equal 35. The accompanying exact checker verifies all $3^6=729$
threshold vectors, with higher thresholds identically zero.

\paragraph{Exact globality certificate.}
Singleton sources have load at most one and never enter the tail.
A subset of the four-element profiles corresponds to a graph on six
vertices, its edges being the omitted pairs. If it has $k$ edges, degrees
$d_i$, and $h\in\{0,1\}$ says whether the three identical all-one sources
are included, then
\[
 D_i=2(k-d_i)+3h,\qquad |A|=k+3h.\tag{6}
\]
The $2^{15}$ graphs have 102 sorted degree sequences. For their 204
cases the supplied JSON gives rational feasible primal covers. The
verifier checks the actual target columns and capacities, coverage of
all degree types, every demand inequality and every exact cost.
All nonempty cases except the full 18-source tail have certified surplus
at most $-7/10$; the empty case gives zero. Thus (5) is the global maximum.
Every pair target is strictly fractionally selected in the displayed
optimal primal, so any optimal dual for this tail obeys $w_i+w_j=1$
for all $i\ne j$. This forces the unique global maximizer $w_i=1/2$.

\paragraph{Every formal cold quotient is negative.}
A nonzero cold formal divisor is exactly one $e_i$. Conditioning on it
and deleting one occurrence leaves 25 sources and 25 targets. After
normalization the prices equal one. The source loads are three fives,
ten sevens and twelve zeros; target loads are six elevens and nineteen
ones. Hence every formal quotient has
\[
 Q_i=3\left(5-\frac{21}{10}\right)+10\left(7-\frac{21}{10}\right)
                    -6(11-1)=-\frac{23}{10}.\tag{7}
\]
Replicating all profiles $L$ times preserves the unique optimum, gives
$J=(6/5)L$, and leaves every cold quotient negative. Thus the abstract
analogue of (D) with error $C\sqrt{90L}$ fails for any fixed $C$.

\paragraph{The precise arithmetic law missing from this model.}
Each formal coordinate occurs in $25L$ of the $90L$ sources. An actual
prime would require $\lfloor90L/p\rfloor=25L$, which is impossible:
for $p\le3$ the count is at least $30L$, and for $p\ge5$ it is at most
$18L$. The violation is already linear in $L$. This model does not
refute arithmetic (D); it proves that abstract full-tail globality,
exact rows and orthant domination alone cannot establish it.

\section{Other checked consequences and the exact remaining gap}
At a joint optimum, work with actual bounded prime-adic phases. For an
active $p$ put $E=\lfloor\log_p M\rfloor+1$ and
$Y_i=\sum_{q\ne p}w_qv_q(x+i)$. CRT permits changing the $p$-phase while
keeping every other prime phase. Minimality of the target hinge and
averaging the $p^E$ phases gives the valid inequality
\[
 H-H_{\setminus p}\le\sum_{j=1}^{E}\frac1{p^j}\sum_{i=1}^{M}
 \big[(Y_i+jw_p-1)_+-(Y_i+(j-1)w_p-1)_+\big].\tag{8}
\]
Its proof expands the hinge in valuation levels and uses probability
$p^{-j}$ for reaching level $j$. It is an upper bound for the target
marginal. Differentiating it as if it were an equality cannot produce
the missing lower payment estimate.

For a fully selected target set $B$, let
$k_{p,j}=\#\{b\in B:p^j\mid b\}$. The difference-product argument gives
\[
 \sum_{p,j}\binom{k_{p,j}}2\log p
 =\sum_{b<b'}\log\gcd(b,b')
 \le\binom{|B|}{2}\log(M-1).
\]
Whole-row balances control sums over $j$, not individual layers.
Increasing $|B|$ weakens the last inequality, so it does not supply the
upper bound on selected cost sought here without an additional argument.

\paragraph{Outcome.}
No proof of (D), no direct $\Delta_c(M)=O_c(\sqrt M)$ bound, and no
counterexample to (D) on genuine arithmetic positive maxima was obtained.
The quotient identity and its conditional closing chain remain valid.
The failed intermediate steps are source-only localization, separate
small-root union transfer, and derivation from abstract optimization
and row data alone.

The unresolved step is a joint arithmetic estimate of $S_c-H$, using
the actual counts $\lfloor M/d\rfloor$ and actual intersections, while
retaining the true source and quotient reserves. The attached verifiers
check the finite certificates used here; they do not compute $g(n)$ or
formally verify a general solution. Detailed working proofs are included
in the same bundle.
\end{document}
